\section{Full Configurations and Architecture Details}
\label{app:configs}
The source package contains all 15 effective config snapshots under
\path{evidence/configs/}. Within each model, seeds change but architecture,
data/evaluation settings and training policy remain fixed. The common protocol
uses Fashion-MNIST, split seed 36, batch 64, Adam at \(10^{-3}\), weight decay
0, a 100-epoch cap and validation macro-F1 checkpoint selection.

\begin{table}[H]\centering\small
\caption{Effective model/training settings; detailed snapshots are included in the source.}
\begin{tabularx}{\linewidth}{lYY}
\toprule Model & Architecture & Model-specific training/regularization\\\midrule
Linear & 784 to 10 & Cosine to \(10^{-5}\); patience 15, min-delta 0.001\\
MLP & 784--256--128--10; ReLU & Hidden dropout 0.1; same cosine/patience as Linear\\
CNN & Channels 32/64; 3-by-3 kernels; max pooling; 64--128--10 head & Head dropout 0.1; no scheduler/early stop\\
GRU & Rows; input 28; hidden 128; one layer, one direction & No inter-layer dropout; no scheduler/early stop\\
Transformer & Patch 4; width 64; four heads; two blocks; FFN width 256 & Dropout 0.1; no scheduler/early stop\\
\bottomrule
\end{tabularx}
\end{table}
The saved configs also identify operational paths and the source configuration
modules. Device/path overrides used for later CPU stages are recorded separately,
not silently presented as the original T4 training setup. Parameter elements are
counted once through the model's parameter iterator; buffers are excluded and
\texttt{requires\_grad} determines trainability.
\FloatBarrier

\section{Per-Seed Results and Run Provenance}
\label{app:runs}
Table~\ref{tab:seeds} reports all seeds. Accuracy/F1 are percentages and the
selected epoch is 1-based; the source history itself uses 0-based indices.
There is no best-test-seed selection.
\begin{table}[H]\centering\footnotesize
\caption{All 15 selected checkpoints and their evaluation results.}
\label{tab:seeds}
\begin{tabular}{lrrrrrrr}
\toprule
Model & Seed & Epochs & Best & Val. acc. & Test acc. & Test F1 & Test CE \\
\midrule
Linear & 69420 & 38 & 30 & 79.16 & 78.12 & 78.098 & 0.61253 \\
Linear & 67 & 33 & 18 & 79.17 & 78.33 & 78.478 & 0.60811 \\
Linear & 69 & 31 & 16 & 79.47 & 78.68 & 78.706 & 0.61362 \\
MLP & 69420 & 100 & 92 & 90.15 & 89.67 & 89.647 & 0.29002 \\
MLP & 67 & 83 & 80 & 90.38 & 89.84 & 89.806 & 0.28514 \\
MLP & 69 & 100 & 93 & 90.47 & 89.65 & 89.641 & 0.28903 \\
CNN & 69420 & 100 & 100 & 91.43 & 90.32 & 90.423 & 0.27405 \\
CNN & 67 & 100 & 100 & 91.66 & 90.92 & 90.885 & 0.25862 \\
CNN & 69 & 100 & 81 & 91.55 & 90.83 & 90.870 & 0.25997 \\
GRU & 69420 & 100 & 62 & 91.13 & 90.38 & 90.366 & 0.28147 \\
GRU & 67 & 100 & 84 & 91.01 & 90.50 & 90.478 & 0.28611 \\
GRU & 69 & 100 & 53 & 91.05 & 90.01 & 90.000 & 0.28312 \\
Transformer & 69420 & 100 & 96 & 91.57 & 91.13 & 91.120 & 0.24628 \\
Transformer & 67 & 100 & 98 & 91.61 & 91.54 & 91.482 & 0.23978 \\
Transformer & 69 & 100 & 100 & 91.56 & 91.13 & 91.100 & 0.24713 \\
\bottomrule
\end{tabular}

\end{table}
The following exact run IDs link these rows to their histories, effective configs
and best checkpoints. Their timestamp and random suffix are artifact identifiers,
not additional experimental factors.
\begin{table}[H]\centering
{\fontsize{9}{11}\selectfont
\setlength{\tabcolsep}{3pt}
\begin{tabular}{lrl}
\toprule
Model & Seed & Run ID \\
\midrule
Linear & 69420 & \texttt{a1\_linear\_main\_seed69420\_20261023-002758\_9761a8} \\
Linear & 67 & \texttt{a1\_linear\_main\_seed67\_20261023-003924\_612a81} \\
Linear & 69 & \texttt{a1\_linear\_main\_seed69\_20261023-004922\_8fcd80} \\
MLP & 69420 & \texttt{a1\_mlp\_main\_seed69420\_20261023-005847\_49b2dc} \\
MLP & 67 & \texttt{a1\_mlp\_main\_seed67\_20261023-013028\_2a1c4c} \\
MLP & 69 & \texttt{a1\_mlp\_main\_seed69\_20261023-015641\_53b249} \\
CNN & 69420 & \texttt{a1\_cnn\_main\_seed69420\_20261023-022827\_cd3858} \\
CNN & 67 & \texttt{a1\_cnn\_main\_seed67\_20261023-030133\_702c4e} \\
CNN & 69 & \texttt{a1\_cnn\_main\_seed69\_20261023-064233\_c73732} \\
GRU & 69420 & \texttt{a1\_rnn\_main\_seed69420\_20261023-071609\_51281c} \\
GRU & 67 & \texttt{a1\_rnn\_main\_seed67\_20261023-074850\_b7dcdb} \\
GRU & 69 & \texttt{a1\_rnn\_main\_seed69\_20261023-082158\_329a4e} \\
Transformer & 69420 & \texttt{a1\_transformer\_main\_seed69420\_20261023-085501\_3eb003} \\
Transformer & 67 & \texttt{a1\_transformer\_main\_seed67\_20261024-030332\_1727fa} \\
Transformer & 69 & \texttt{a1\_transformer\_main\_seed69\_20261024-034117\_45d824} \\
\bottomrule
\end{tabular}
}
\caption{Exact run identifiers. The \texttt{rnn} artifact name denotes the GRU configuration.}
\label{tab:ids}
\end{table}
The source commit and common identities recorded by the bundle are:
\begin{description}
\item[Git revision] \small\texttt{a8539bee778d99e2902a1dd36aa157c6ebb3a42c}
\item[Split SHA-256] \small\texttt{e4a776305b2a1b5c392ba196d023de004f47124e8f831495b9733c48e72c25e6}
\item[Statistics SHA-256] \small\texttt{e47b863cb866cdc61613bb88d772d02f8aedc0a21fa56b79dee27b27d341381e}
\end{description}
Full per-run checkpoint/config hashes are retained in
\path{evidence/report_facts.json} and the original run metadata.
\FloatBarrier

\section{Additional Class-Level Results}
\label{app:class}
The tables below aggregate official-test precision, recall and F1 over the three
seeds as mean \(\pm\) sample SD, expressed in percent. Support is 1,000 examples
per class per run. The same test images are evaluated repeatedly; the support
is not 3,000 independent images. Validation and per-seed class results are also
included in the machine-readable \path{evidence/per_class.csv}.

\begin{table}[H]\centering\small
\caption{Linear official-test per-class metrics, mean \(\pm\) sample SD, \(n=3\).}
\begin{tabular}{lrrrr}
\toprule
Class & Precision (\%) & Recall (\%) & F1 (\%) & Support \\
\midrule
T-shirt/top & 77.12 $\pm$ 1.72 & 76.90 $\pm$ 1.35 & 76.99 $\pm$ 0.26 & 1000 \\
Trouser & 89.83 $\pm$ 3.06 & 88.57 $\pm$ 2.74 & 89.15 $\pm$ 1.54 & 1000 \\
Pullover & 60.56 $\pm$ 3.01 & 66.00 $\pm$ 4.35 & 63.07 $\pm$ 2.20 & 1000 \\
Dress & 75.08 $\pm$ 1.77 & 79.03 $\pm$ 2.55 & 76.99 $\pm$ 1.51 & 1000 \\
Coat & 71.49 $\pm$ 1.08 & 56.17 $\pm$ 5.76 & 62.80 $\pm$ 3.74 & 1000 \\
Sandal & 92.84 $\pm$ 1.18 & 88.73 $\pm$ 1.42 & 90.73 $\pm$ 0.21 & 1000 \\
Shirt & 47.41 $\pm$ 0.79 & 51.80 $\pm$ 1.30 & 49.50 $\pm$ 0.28 & 1000 \\
Sneaker & 86.42 $\pm$ 2.05 & 93.53 $\pm$ 1.53 & 89.81 $\pm$ 0.42 & 1000 \\
Bag & 93.28 $\pm$ 0.54 & 91.50 $\pm$ 0.44 & 92.38 $\pm$ 0.11 & 1000 \\
Ankle boot & 94.28 $\pm$ 1.69 & 91.53 $\pm$ 1.98 & 92.86 $\pm$ 0.23 & 1000 \\
\bottomrule
\end{tabular}

\end{table}

\begin{table}[H]\centering\small
\caption{MLP official-test per-class metrics, mean \(\pm\) sample SD, \(n=3\).}
\begin{tabular}{lrrrr}
\toprule
Class & Precision (\%) & Recall (\%) & F1 (\%) & Support \\
\midrule
T-shirt/top & 84.73 $\pm$ 0.22 & 84.53 $\pm$ 0.47 & 84.63 $\pm$ 0.32 & 1000 \\
Trouser & 99.22 $\pm$ 0.38 & 97.67 $\pm$ 0.38 & 98.44 $\pm$ 0.33 & 1000 \\
Pullover & 84.06 $\pm$ 0.98 & 82.67 $\pm$ 1.40 & 83.35 $\pm$ 0.24 & 1000 \\
Dress & 88.89 $\pm$ 0.33 & 90.67 $\pm$ 0.15 & 89.77 $\pm$ 0.19 & 1000 \\
Coat & 80.11 $\pm$ 0.26 & 84.20 $\pm$ 0.40 & 82.11 $\pm$ 0.23 & 1000 \\
Sandal & 97.62 $\pm$ 0.29 & 96.93 $\pm$ 0.75 & 97.27 $\pm$ 0.31 & 1000 \\
Shirt & 73.95 $\pm$ 0.95 & 70.53 $\pm$ 1.01 & 72.19 $\pm$ 0.43 & 1000 \\
Sneaker & 94.18 $\pm$ 0.06 & 96.63 $\pm$ 0.12 & 95.39 $\pm$ 0.08 & 1000 \\
Bag & 97.30 $\pm$ 0.10 & 97.43 $\pm$ 0.06 & 97.37 $\pm$ 0.03 & 1000 \\
Ankle boot & 97.00 $\pm$ 0.49 & 95.93 $\pm$ 0.38 & 96.46 $\pm$ 0.17 & 1000 \\
\bottomrule
\end{tabular}

\end{table}

\begin{table}[H]\centering\small
\caption{CNN official-test per-class metrics, mean \(\pm\) sample SD, \(n=3\).}
\begin{tabular}{lrrrr}
\toprule
Class & Precision (\%) & Recall (\%) & F1 (\%) & Support \\
\midrule
T-shirt/top & 86.32 $\pm$ 2.79 & 84.27 $\pm$ 3.42 & 85.21 $\pm$ 0.47 & 1000 \\
Trouser & 99.26 $\pm$ 0.21 & 97.57 $\pm$ 0.59 & 98.40 $\pm$ 0.20 & 1000 \\
Pullover & 87.28 $\pm$ 1.26 & 87.20 $\pm$ 1.04 & 87.23 $\pm$ 0.16 & 1000 \\
Dress & 89.51 $\pm$ 1.05 & 89.20 $\pm$ 1.06 & 89.35 $\pm$ 0.31 & 1000 \\
Coat & 85.34 $\pm$ 0.25 & 85.13 $\pm$ 1.40 & 85.23 $\pm$ 0.61 & 1000 \\
Sandal & 98.18 $\pm$ 0.44 & 96.63 $\pm$ 1.59 & 97.39 $\pm$ 0.69 & 1000 \\
Shirt & 73.15 $\pm$ 3.82 & 76.13 $\pm$ 4.20 & 74.47 $\pm$ 0.42 & 1000 \\
Sneaker & 94.17 $\pm$ 2.03 & 96.57 $\pm$ 1.52 & 95.33 $\pm$ 0.36 & 1000 \\
Bag & 98.36 $\pm$ 0.24 & 98.20 $\pm$ 0.35 & 98.28 $\pm$ 0.11 & 1000 \\
Ankle boot & 96.75 $\pm$ 0.88 & 96.00 $\pm$ 1.08 & 96.37 $\pm$ 0.11 & 1000 \\
\bottomrule
\end{tabular}

\end{table}

\begin{table}[H]\centering\small
\caption{GRU official-test per-class metrics, mean \(\pm\) sample SD, \(n=3\).}
\begin{tabular}{lrrrr}
\toprule
Class & Precision (\%) & Recall (\%) & F1 (\%) & Support \\
\midrule
T-shirt/top & 85.39 $\pm$ 0.87 & 82.33 $\pm$ 1.50 & 83.82 $\pm$ 0.43 & 1000 \\
Trouser & 99.06 $\pm$ 0.40 & 98.13 $\pm$ 0.32 & 98.59 $\pm$ 0.10 & 1000 \\
Pullover & 84.19 $\pm$ 0.98 & 83.73 $\pm$ 0.21 & 83.96 $\pm$ 0.58 & 1000 \\
Dress & 88.91 $\pm$ 0.41 & 91.33 $\pm$ 1.47 & 90.10 $\pm$ 0.62 & 1000 \\
Coat & 83.43 $\pm$ 1.10 & 84.60 $\pm$ 1.85 & 83.99 $\pm$ 0.36 & 1000 \\
Sandal & 97.76 $\pm$ 0.31 & 97.27 $\pm$ 0.40 & 97.51 $\pm$ 0.34 & 1000 \\
Shirt & 74.46 $\pm$ 1.34 & 73.97 $\pm$ 1.29 & 74.21 $\pm$ 1.31 & 1000 \\
Sneaker & 94.77 $\pm$ 0.81 & 97.23 $\pm$ 0.32 & 95.99 $\pm$ 0.56 & 1000 \\
Bag & 97.21 $\pm$ 0.77 & 98.67 $\pm$ 0.12 & 97.93 $\pm$ 0.34 & 1000 \\
Ankle boot & 97.72 $\pm$ 0.07 & 95.70 $\pm$ 0.66 & 96.70 $\pm$ 0.37 & 1000 \\
\bottomrule
\end{tabular}

\end{table}

\begin{table}[H]\centering\small
\caption{Transformer official-test per-class metrics, mean \(\pm\) sample SD, \(n=3\).}
\begin{tabular}{lrrrr}
\toprule
Class & Precision (\%) & Recall (\%) & F1 (\%) & Support \\
\midrule
T-shirt/top & 87.10 $\pm$ 0.55 & 85.50 $\pm$ 0.17 & 86.29 $\pm$ 0.18 & 1000 \\
Trouser & 98.73 $\pm$ 0.54 & 98.50 $\pm$ 0.53 & 98.62 $\pm$ 0.33 & 1000 \\
Pullover & 86.37 $\pm$ 1.04 & 88.27 $\pm$ 0.93 & 87.31 $\pm$ 0.79 & 1000 \\
Dress & 91.04 $\pm$ 2.87 & 91.27 $\pm$ 2.00 & 91.11 $\pm$ 0.54 & 1000 \\
Coat & 83.06 $\pm$ 1.48 & 87.10 $\pm$ 2.95 & 84.99 $\pm$ 0.67 & 1000 \\
Sandal & 98.20 $\pm$ 0.17 & 98.00 $\pm$ 0.17 & 98.10 $\pm$ 0.17 & 1000 \\
Shirt & 77.20 $\pm$ 1.44 & 72.67 $\pm$ 1.37 & 74.85 $\pm$ 0.49 & 1000 \\
Sneaker & 95.58 $\pm$ 0.58 & 96.47 $\pm$ 0.84 & 96.02 $\pm$ 0.37 & 1000 \\
Bag & 98.34 $\pm$ 0.06 & 98.47 $\pm$ 0.15 & 98.40 $\pm$ 0.10 & 1000 \\
Ankle boot & 96.89 $\pm$ 0.69 & 96.43 $\pm$ 0.38 & 96.66 $\pm$ 0.22 & 1000 \\
\bottomrule
\end{tabular}

\end{table}

\FloatBarrier
\section{Detailed Verification and Reproduction Notes}
\label{app:verification}
The saved-output audit of 24 October 2026 verified 382 files and 61 completed
stages in the original CPU evaluation profile, independently recomputing metrics
from 300,000 prediction rows across 30 evaluations. Class order, argmax labels,
support, confusion counts, gallery selection and sample IDs agree. Mean negative
log true-class probability matches the saved cross-entropy within \(10^{-6}\).
The replay agreement supports using the CPU evaluations of the T4-selected
checkpoints; it is not evidence of bitwise equivalence on every platform.

The retained Windows timing profile, \texttt{cpu-f7208f3550fc}, contains all
15 checkpoint identities and 100 raw batch timings per run. Its session and
per-run records are copied into the source package. The original CPU evaluation
profile, \texttt{cpu-0bd6ba8e08e9}, explicitly skipped GPU timing. A separate
Colab timing summary appeared in an intermediate report without its raw session
records; this final document uses the retained Windows measurements instead.
No new inference or training was run to construct the final tables.

Report-generation checks on 26 October compare all classification aggregates
with the retained evaluation summary, verify checkpoint hashes, inspect official
IDX/gzip identities and class counts, regenerate figures from saved outputs,
and check LaTeX references and layout. The example preprocessing batch is newly
regenerated data, clearly separated from original training history.

Source line endings deserve care: the committed split JSON and the Windows
checkout can have different byte hashes because of LF/CRLF conversion while
parsing to the same indices. Provenance comparisons use the recorded committed
blob identity and semantic membership, rather than silently relabeling an
unrelated checkout-file hash.

The editable content is organized under \texttt{sections/}, \texttt{tables/}
and \texttt{figures/}. The \texttt{evidence/} directory retains effective
configs, per-run/class aggregates, raw Windows timing samples, resource hashes
and figure selections. The \texttt{reproduction/} directory provides the
pinned-source reconstruction helper and standalone Colab evaluation/timing
notebooks. Together, these records distinguish original measurements from
later report-generation checks.
